\section{Retained proofs for operational and vector certificates}\label{supp:relocated69}
\subsection{Proof retained from operational67}\label{supp:relocated69-1}
\begin{proof}
Strong convexity gives $f((1-s)l+su)\le(1-s)f(l)+sf(u)-\mu s(1-s)(u-l)^2/2$. Replacing the endpoint values by their upper bounds proves \eqref{eq:screen67}. Subtraction of a valid lower bound proves \eqref{eq:screen-stop67}. A failed screen adds no endpoint to the partition and alters no lower bound. Subsequent protected refinement is therefore covered by the earlier termination proof. A successful screen supplies exactly the local accuracy premise used in the acquired-policy recursion. Finally, $C$ is a convex quadratic in $a$, proving the sublevel-set assertion.
\end{proof}

\subsection{Proof retained from operational67}\label{supp:relocated69-2}
\begin{proof}
A successful screen returns without an endpoint query. A failed screen invokes a procedure bounded by $K(z)$, and its margin loss equals one. This proves \eqref{eq:work-bound67}. For each fixed predictor the variable $K(z)\ell_\tau(z,\theta)$ lies in $[0,K_{\max}]$. The bounded-sum inequality of \citet{hoeffding1963} and a union bound over the conditionally fixed family give \eqref{eq:work-validation67}.
\end{proof}

\subsection{Proof retained from operational67}\label{supp:relocated69-3}
\begin{proof}
The chord is a convex quadratic in its normalized action, which is affine in $w$. Squaring the positive part preserves convexity and gives a continuously differentiable function. The nonnegative weighted sum and quadratic penalty preserve these properties. Convexity gives $F(w)-F(v)\le g^\top(w-v)$, and the minimum of $g^\top v$ on the box is $\sum_j\min\{0,g_j\}$. This proves \eqref{eq:dual-gap67}. Every failed screen has squared positive-part residual at least $\tau^2$, while every successful screen uses zero further endpoint queries. Summing the corresponding $K_i$ bounds proves the work claim. Strong convexity with modulus $\lambda$ proves the distance bound.
\end{proof}

\subsection{Proof retained from vector67}\label{supp:relocated69-4}
\begin{proof}
Semiconcavity means that $f(a)-H\|a\|_2^2/2$ is concave. Its barycentric inequality, with vertex values replaced by lower endpoints, gives \eqref{eq:triangle67}. The tangent inequality for the convex quadratic $p$ proves \eqref{eq:dual-triangle67}; a linear function attains its simplex minimum at a vertex. Taking minima over a cover preserves the lower bound. Each feasible upper witness is an action of the actual model. The local gap premise therefore implies the earlier policy recursion, irrespective of how the triangle or support point was proposed.

For finiteness, use simplex lattice denominator $A$, with $(A+1)(A+2)/2$ vertices. On triangles of diameter proportional to $c(x)/A$, the curvature error tends to zero as $A^{-2}$. The endpoint error is reduced by finite descendant accuracy and quadrature refinements. The retained exact-rational recovery contract requires $\delta>16\Lambda\,2^{-32}$ at each recursive call, leaving positive slack after rounding; dyadic enlargement of $A$ and $q$ then terminates. This is a finite arithmetic procedure, not a uniform bound on rational bit cost. For \eqref{eq:two-shock67}, the independent innovation variances after conditioning on the bins are $\rho_j^2/(3q^2)$, their direction norms squared are $d$, and the semiconcavity Taylor remainder is $M_{t+1}/2$ times their sum. Substitution of $\rho_1=1/32$, $\rho_2=1/64$ gives the coefficient. Finally, each action query has $q_r^2$ descendants, proving \eqref{eq:vector-work67}.
\end{proof}
